% Figure: catastrophic cancellation. Subtracting two numbers that agree
% to k significant figures amplifies relative rounding error by about
% 10^k. The plotted line is the predicted relative-error floor
% (unit roundoff) times the condition number, in double precision.
\begin{figure}[htbp]
\centering
\begin{tikzpicture}
\begin{semilogyaxis}[
  width=11cm, height=6.5cm,
  xlabel={agreeing significant figures $k$ of the two operands},
  ylabel={relative error of the difference},
  xmin=0, xmax=15,
  ymin=1e-16, ymax=1e0,
  grid=both,
  grid style={engblack!12},
  tick label style={font=\footnotesize},
  label style={font=\small},
  legend style={font=\footnotesize, at={(0.02,0.98)}, anchor=north west},
]
% predicted floor: u * 10^k, with u = 1.1e-16 (double precision)
\addplot[engblue, very thick, domain=0:15, samples=16]
  {1.1e-16 * 10^(x)};
\addlegendentry{$u \cdot 10^{k}$ (double precision)}
% the unit-roundoff floor for reference
\addplot[engblack!50, dashed, domain=0:15] {1.1e-16};
\addlegendentry{unit roundoff $u$}
% a marked threshold where half the digits are gone
\addplot[only marks, mark=*, engred, mark size=2.2pt]
  coordinates {(8, 1.1e-8)};
\node[font=\footnotesize, anchor=west, engred]
  at (axis cs:8.3, 1.1e-8) {8 figures lost};
\end{semilogyaxis}
\end{tikzpicture}
\caption[Catastrophic cancellation amplifies relative error]{Subtracting
two positive numbers that agree to $k$ significant figures multiplies
the relative rounding error of the operands by about $10^{k}$. Each
operand starts at the unit-roundoff floor $u \approx 1.1\times10^{-16}$
of IEEE 754 double precision (dashed line); the difference carries the
amplified relative error $u\cdot 10^{k}$ (solid line). When the operands
agree to eight figures, the difference has surrendered eight of the
sixteen figures double precision held. The remedy is an algebraic
rewrite, such as the Vieta companion for a quadratic root or the
conjugate multiplication for $\sqrt{x+1}-\sqrt{x}$, that removes the
cancelling subtraction.}
\label{fig:vol02:ch01:cancellation}
\end{figure}
